%! Solving Radical Equations \& Extraneous Solutions — Unit 5, Lesson 5.5 — Slides (Beamer)
\documentclass[aspectratio=169, 11pt]{beamer}
\usepackage{algebra2-beamer}   % provides \forestheader{} and \sectionlabel[color]{}
% Standard coordinate grid + axes: {xmin}{xmax}{ymin}{ymax}
\newcommand{\lgrid}[4]{%
  \draw[step=1, linegray!45, thin] (#1,#3) grid (#2,#4);
  \draw[->, thick, charcoal] (#1,0)--(#2,0) node[below right, font=\tiny]{$x$};
  \draw[->, thick, charcoal] (0,#3)--(0,#4) node[left, font=\tiny]{$y$};}

\begin{document}

% TITLE SLIDE (CourseName is not defined in beamer, so the name is literal)
{
\setbeamercolor{background canvas}{bg=forest}
\begin{frame}[plain]
  \vspace{2.8cm}\hspace{0.55cm}%
  \begin{minipage}{0.9\textwidth}
    {\color{goldacc}\bfseries\small LESSON 5.5}\\[0.5cm]
    {\color{white}\bfseries\LARGE Solving Radical Equations \& Extraneous Solutions}\\[1.2cm]
    {\color{forestmid}\small Algebra 2~$\cdot$~Unit 5: Radical Functions}
  \end{minipage}
\end{frame}
}

% HOOK
\begin{frame}[t, plain]
  \forestheader{Every line is correct. The answer is not.}
  \sectionlabel{Hook}
  \begin{columns}[T]
    \begin{column}{0.44\textwidth}
      \footnotesize
      Solve $\sqrt{x+7}=x-5$:
      \vspace{3pt}
      \begin{center}
      \renewcommand{\arraystretch}{1.3}
      \begin{tabular}{@{}l l@{}}
      $\sqrt{x+7}=x-5$ & original \\
      $x+7=(x-5)^{2}$ & square \\
      $x+7=x^{2}-10x+25$ & expand \\
      $0=x^{2}-11x+18$ & collect \\
      $0=(x-2)(x-9)$ & factor \\
      $x=2$ \; or \; $x=9$ & \\
      \end{tabular}
      \end{center}
      \vspace{2pt}
      Find the bad step. \textbf{There isn't one.}
    \end{column}
    \begin{column}{0.52\textwidth}
      \footnotesize
      Now test both in the \emph{original}:
      \vspace{3pt}
      \begin{center}
      \renewcommand{\arraystretch}{1.3}
      \begin{tabular}{|c|c|c|c|}
      \hline
       & $\sqrt{x+7}$ & $x-5$ & Equal? \\ \hline
      $x=9$ & $4$ & $4$ & yes \\ \hline
      $x=2$ & $3$ & $-3$ & \textbf{no} \\ \hline
      \end{tabular}
      \end{center}
      \vspace{4pt}
      \begin{block}{The question the whole lesson answers}
        One answer is wrong, and \emph{no step} was wrong. \\[3pt]
        So where did the wrong answer come from?
      \end{block}
    \end{column}
  \end{columns}
  \vspace{2pt}
  \begin{center}
    \textbf{Squaring hands you candidates, not solutions. The check is part of the method.}
  \end{center}
\end{frame}

% THE METHOD
\begin{frame}[t, plain]
  \forestheader{Isolate, raise, solve, check}
  \sectionlabel[goldacc]{The method}
  \begin{columns}[T]
    \begin{column}{0.50\textwidth}
      \footnotesize
      \begin{enumerate}\setlength{\itemsep}{4pt}
        \item \textbf{Isolate} the radical --- alone on one side.
        \item Raise \emph{both sides} to the power equal to the \textbf{index}.
        \item Solve what is left (linear or quadratic).
        \item \textbf{Check every answer in the original equation.}
      \end{enumerate}
      \vspace{4pt}
      \begin{block}{Why Step 1 is first}
        $\left(\sqrt{x-3}\right)^{2}=x-3$ \quad --- clean. \\[3pt]
        $\left(\sqrt{x}+3\right)^{2}=x+6\sqrt{x}+9$ \quad --- \emph{worse}. \\[3pt]
        Squaring erases a radical only when it stands \textbf{alone}.
      \end{block}
    \end{column}
    \begin{column}{0.46\textwidth}
      \footnotesize
      \textbf{Worked:} $\sqrt{x+4}-1=3$
      \vspace{3pt}
      \begin{center}
      \renewcommand{\arraystretch}{1.4}
      \begin{tabular}{@{}l l@{}}
      $\sqrt{x+4}=4$ & isolate \\
      $x+4=16$ & square \\
      $x=12$ & solve \\
      \end{tabular}
      \end{center}
      \vspace{3pt}
      Check: $\sqrt{16}-1=4-1=3$ \checkmark
      \vspace{4pt}
      \begin{block}{What the equation is really asking}
        This is 5.0's $f(x)=\sqrt{x+4}-1$. \\[3pt]
        \emph{Where does that graph equal $3$?}
      \end{block}
    \end{column}
  \end{columns}
\end{frame}

% THE ONE-WAY STREET
\begin{frame}[t, plain]
  \forestheader{Squaring is a one-way street}
  \sectionlabel[goldacc]{Why the check exists}
  \begin{columns}[T]
    \begin{column}{0.46\textwidth}
      \footnotesize
      \begin{block}{The whole mechanism}
        If $a=b$, then $a^{2}=b^{2}$ --- \emph{always}. \\[4pt]
        But if $a^{2}=b^{2}$, then $a=b$ \quad\textbf{or}\quad $a=-b$.
      \end{block}
      \vspace{3pt}
      \begin{itemize}\setlength{\itemsep}{3pt}
        \item Squaring can never \emph{lose} a solution.
        \item Squaring can always \emph{gain} one.
        \item $\sqrt9$ and $-3$ are different numbers --- but their squares are the same number.
      \end{itemize}
      \vspace{3pt}
      \begin{center}
      \textbf{Squaring destroys sign information.}
      \end{center}
    \end{column}
    \begin{column}{0.50\textwidth}
      \centering
      \begin{tikzpicture}[scale=0.30]
        \lgrid{-8}{13}{-6}{6}
        \begin{scope}
          \clip (-8,-6) rectangle (13,6);
          \draw[very thick, forest, domain=-7:13, samples=110, smooth] plot (\x,{sqrt((\x)+7)});
          \draw[very thick, navy, domain=-1:11, samples=2] plot (\x,{(\x)-5});
        \end{scope}
        \draw[dashed, redacc, very thick] (2,3)--(2,-3);
        \fill[forest] (-7,0) circle (6pt);
        \fill[charcoal] (9,4) circle (5pt);
        \fill[redacc] (2,3) circle (5pt);
        \fill[redacc] (2,-3) circle (5pt);
      \end{tikzpicture}\\[2pt]
      {\scriptsize \textcolor{forest}{$y=\sqrt{x+7}$} \; and \; \textcolor{navy}{$y=x-5$} --- they
      cross \textbf{once}, at $(9,4)$. \\
      At $x=2$: the curve is at $3$, the line at $-3$. Both square to $9$.}
    \end{column}
  \end{columns}
\end{frame}

% THE SHORTCUT
\begin{frame}[t, plain]
  \forestheader{Read the equation before you solve it}
  \sectionlabel[goldacc]{A free shortcut}
  \begin{columns}[T]
    \begin{column}{0.48\textwidth}
      \footnotesize
      \textbf{Solve} $\sqrt{2x-1}+5=2$.
      \vspace{4pt}
      \[ \sqrt{2x-1}=-3 \]
      \vspace{-6pt}
      \begin{block}{Stop here}
        From Lesson 5.1: an even root is the \textbf{principal} root, and a principal root is
        \emph{never negative}. \\[3pt]
        \textbf{No real solution} --- and you wrote one line.
      \end{block}
    \end{column}
    \begin{column}{0.48\textwidth}
      \footnotesize
      \textbf{Square it anyway and watch:}
      \vspace{3pt}
      \[ 2x-1=9 \;\Rightarrow\; x=5 \]
      Check: $\sqrt{2(5)-1}+5=3+5=8 \ne 2$. \\[4pt]
      \textbf{Extraneous.}
      \vspace{4pt}
      \begin{block}{The point}
        The algebra never lied to you. \\[3pt]
        It just never \emph{warned} you.
      \end{block}
    \end{column}
  \end{columns}
\end{frame}

% ODD INDEX
\begin{frame}[t, plain]
  \forestheader{Change the index and the trapdoor closes}
  \sectionlabel{Even vs.\ odd}
  \begin{columns}[T]
    \begin{column}{0.48\textwidth}
      \footnotesize
      \begin{center}
      \renewcommand{\arraystretch}{1.4}
      \begin{tabular}{|l|l|}
      \hline
      \textbf{$\sqrt{x-4}=-2$} & \textbf{$\sqrt[3]{x-4}=-2$} \\ \hline
      no real solution & $x-4=-8$ \\
      & $x=-4$ \\
      & Check: $\sqrt[3]{-8}=-2$ \checkmark \\ \hline
      \end{tabular}
      \end{center}
      \vspace{4pt}
      \begin{center}
      Same numbers. Opposite verdicts. \\
      \textbf{The only difference is the index.}
      \end{center}
    \end{column}
    \begin{column}{0.48\textwidth}
      \footnotesize
      \begin{block}{Why an odd index is safe}
        Every real number has exactly \textbf{one} real cube root. \\[3pt]
        So $A^{3}=B^{3}$ forces $A=B$ --- cubing is \textbf{reversible}, and a reversible step
        cannot invent solutions.
      \end{block}
      \vspace{3pt}
      \textbf{Worked:} $2\sqrt[3]{x+1}-3=1$
      \[ \sqrt[3]{x+1}=2 \;\Rightarrow\; x+1=8 \;\Rightarrow\; x=7 \]
      \emph{Still check.} An odd index protects you from phantoms, not from arithmetic.
    \end{column}
  \end{columns}
\end{frame}

% GRAPHICAL
\begin{frame}[t, plain]
  \forestheader{How many answers should you expect?}
  \sectionlabel[goldacc]{Verify on a graph}
  \begin{columns}[T]
    \begin{column}{0.52\textwidth}
      \centering
      \begin{tikzpicture}[scale=0.30]
        \lgrid{-5}{15}{-3}{5}
        \begin{scope}
          \clip (-5,-3) rectangle (15,5);
          \draw[very thick, forest, domain=-4:15, samples=110, smooth] plot (\x,{sqrt((\x)+4)-1});
          \draw[very thick, navy] (-5,3)--(15,3);
          \draw[very thick, goldacc] (-5,-2)--(15,-2);
        \end{scope}
        \fill[forest] (-4,-1) circle (6pt);
        \fill[charcoal] (12,3) circle (5pt);
      \end{tikzpicture}\\[2pt]
      {\scriptsize \textcolor{forest}{$y=\sqrt{x+4}-1$} \; with \; \textcolor{navy}{$y=3$} \; and \;
      \textcolor{goldacc}{$y=-2$}}
    \end{column}
    \begin{column}{0.44\textwidth}
      \footnotesize
      \begin{itemize}\setlength{\itemsep}{4pt}
        \item $\sqrt{x+4}-1=3$: one crossing, at $x=12$ --- the same answer the algebra gave.
        \item $\sqrt{x+4}-1=-2$: the line passes \emph{below} the endpoint. No crossing, no
              solution.
      \end{itemize}
      \vspace{3pt}
      \begin{block}{The rule}
        The number of \textbf{intersections} is the number of real solutions. \\[3pt]
        More candidates than crossings $\Rightarrow$ the extras are extraneous.
      \end{block}
    \end{column}
  \end{columns}
\end{frame}

% RATIONAL EXPONENTS
\begin{frame}[t, plain]
  \forestheader{The same method in exponent clothing}
  \sectionlabel[goldacc]{$x^{m/n}=k$}
  \begin{columns}[T]
    \begin{column}{0.48\textwidth}
      \footnotesize
      To undo the power $\frac{m}{n}$, raise both sides to its \textbf{reciprocal} $\frac{n}{m}$,
      because $\left(x^{m/n}\right)^{n/m}=x^{1}$.
      \vspace{4pt}
      \[ x^{3/2}=27 \;\Rightarrow\; x=27^{2/3}=\left(\sqrt[3]{27}\right)^{2}=9 \]
      Check: $9^{3/2}=\left(\sqrt9\right)^{3}=27$ \checkmark
    \end{column}
    \begin{column}{0.48\textwidth}
      \footnotesize
      \[ x^{2/3}=9 \;\Rightarrow\; x=\pm 9^{3/2}=\pm 27 \]
      Both check: $(-27)^{2/3}=(-3)^{2}=9$ \checkmark
      \vspace{4pt}
      \begin{block}{Read the fraction first}
        \textbf{Even numerator} $=$ a square in disguise --- write the $\pm$ yourself. \\[3pt]
        \textbf{Even denominator} $=$ an even root --- the base cannot be negative, so a negative
        answer would be extraneous.
      \end{block}
    \end{column}
  \end{columns}
\end{frame}

% SUMMARY
\begin{frame}[t, plain]
  \forestheader{Candidate or solution?}
  \sectionlabel{Summary}
  \begin{center}
  \renewcommand{\arraystretch}{1.35}
  \begin{tabular}{lll}
  \hline
   & \textbf{Even index ($\sqrt{\ }$, $\sqrt[4]{\ }$)} & \textbf{Odd index ($\sqrt[3]{\ }$)} \\
  \hline
  Can the root be negative? & no --- principal root & yes \\
  Raising both sides is & \emph{not} reversible & reversible \\
  Can it invent solutions? & \textbf{yes} & no \\
  Is the check required? & \textbf{yes} & yes, for arithmetic \\
  Root $=$ negative number & no real solution & solve normally \\
  \hline
  \end{tabular}
  \end{center}
  \vspace{6pt}
  \begin{center}
    A candidate is extraneous exactly when it makes the \textbf{other side negative}. \\[3pt]
    Being a negative number is \emph{not} the same thing --- and never was.
  \end{center}
\end{frame}

% ROADMAP
\begin{frame}[t, plain]
  \forestheader{Where Unit 5 goes next}
  \sectionlabel{Connections}
  \textbf{Today:} isolate, raise to the $n$th power, solve --- and check, because squaring destroys
  sign information and hands back answers the original equation never accepted.\\[8pt]
  \begin{itemize}\setlength{\itemsep}{3pt}
    \item \textbf{4.6} Rational equations --- the same structure: a non-reversible step produces
          candidates, and only the check promotes them.
    \item \textbf{5.1--5.4} Principal roots, the index rule, and the graphs you just used as the two
          sides of an equation.
    \item \textbf{5.6} Composition of functions --- one operation feeding another, $(f\circ
          g)(x)=f(g(x))$, where the order matters.
    \item \textbf{5.7} Inverse functions --- composition becomes the \emph{test}, and 5.0's
          reflection over $y=x$ finally gets its explanation.
  \end{itemize}
\end{frame}

\end{document}
